% GENERATED FILE. Edit canonical lessons, then run npm run generate:books.
% canonical-source-hash: fnv1a64:717feb14774bfd23
% canonical-lessons: MR-C43-javal, MR-C43-samor, MR-C43-laamb, MR-R43-how-far

\chapter{Near, Opposite, Far}
\label{ch:mr-near-opposite-far}
\hlchaptermodality{43}

\begin{chapteropening}
\textbf{By the end of this chapter:} \emph{I can place something near or in front of a landmark, say that it is far, and tell a postposition apart from an ordinary word by whether it bends the noun in front of it.}

Complete MR-R43-how-far and answer one where-question three ways -- near Pune, in front of the room, far -- naming which of the three needed no ending at all.
\end{chapteropening}

\section[javaḷ]{\mr{जवळ} (javaḷ) — near, close to}

\label{lesson:MR-C43-javal}



\begin{quote}
Last chapter stacked things above and below each other. Real directions are almost never that tidy.

\begin{itemize}
\raggedright
  \item \emph{Recall:} say \emph{khālī}, then \emph{var}, then bend \textbf{\mr{घर}} and glue each on
\end{itemize}
\end{quote}

\subsection*{You'll want to know: \mr{जवळ}}
\begin{quote}
\textbf{\mr{जवळ}} — \emph{javaḷ} — \textbf{near, close to}
\end{quote}

\begin{quote}
\textbf{\mr{माझं} \mr{घर} \mr{पुण्याजवळ} \mr{आहे}.} — \emph{mājhaṁ ghar puṇyājavaḷ āhe.} — \textbf{my house is near Pune.}
\end{quote}

This is the answer a real person gives when you ask \textbf{\mr{कुठे}?} and the true answer is not a single named point. It takes the same oblique stem as everything else in this column — \textbf{\mr{पुणे} $\to$ \mr{पुण्या}-} — so nothing new has to be learned about the noun.

\textbf{\mr{जवळ}} carries the \textbf{\mr{ळ}}, the retroflex \emph{l} that is Marathi's own letter and that Hindi does not have. Say it with the tongue curled back, the way you learned it in chapter three, or the word will sound like a different one.

It has a second job you will need in the next chapter: \textbf{X \mr{जवळ}} can also mean \emph{in X's keeping} — near enough to reach — which is how Marathi gets close to saying somebody HAS something.

\subsection*{Your turn}
Say these aloud:

\begin{itemize}
\raggedright
  \item \emph{javaḷ} — near, with the tongue curled on the \textbf{\mr{ळ}}
  \item \emph{mājhaṁ ghar puṇyājavaḷ āhe}
  \item \emph{mitrājavaḷ} — near my friend
\end{itemize}

\begin{tcolorbox}[breakable,colback=teal!4,colframe=teal!35!black,title={Before you move on}]
Which letter in \textbf{\mr{जवळ}} is the one Hindi does not have? (\textbf{\mr{ळ}}.) What does \textbf{\mr{पुणे}} become before it? (\textbf{\mr{पुण्या}-}.)

Sources: \href{https://en.wiktionary.org/wiki/\%E0\%A4\%9C\%E0\%A4\%B5\%E0\%A4\%B3}{Wiktionary: \mr{जवळ}}.
\end{tcolorbox}
\section[samor]{\mr{समोर} (samor) — in front of, opposite}

\label{lesson:MR-C43-samor}



\begin{quote}
\textbf{\mr{जवळ}} says \emph{close by}, which does not say which side.
\end{quote}

\subsection*{You'll want to know: \mr{समोर}}
\begin{quote}
\textbf{\mr{समोर}} — \emph{samor} — \textbf{in front of, opposite, facing}
\end{quote}

\begin{quote}
\textbf{\mr{घरासमोर}} — \emph{gharāsamor} — \textbf{in front of the house}
\end{quote}

This is the postposition a set of directions is actually made of. \emph{Opposite the station}, \emph{in front of the shop} — a Marathi speaker giving you a landmark reaches for \textbf{\mr{समोर}} before anything else.

Take the word apart and it pays you back. It is Sanskrit \textbf{\mr{सम्मुख}} (\emph{sammukha}) — \textbf{\mr{सम्}} \emph{together} plus \textbf{\mr{मुख}} \emph{face}. Chapter nineteen told you that \textbf{\mr{मुख}} is the formal doublet of \textbf{\mr{तोंड}}, the everyday mouth. So \textbf{\mr{समोर}} is, quite literally, \textbf{face-to-face}: the thing whose face is turned toward yours.

That is why it means \emph{opposite} as well as \emph{in front of}. English keeps those apart; Marathi does not need to, because a face has only one direction.

\subsection*{Your turn}
Say these aloud:

\begin{itemize}
\raggedright
  \item \emph{samor} — in front of
  \item \emph{gharāsamor}, then \emph{gharājavaḷ} — a side, then a distance
  \item which body part is inside this word — \emph{\textbf{\mr{मुख}}}, the formal \emph{tond}
\end{itemize}

\begin{tcolorbox}[breakable,colback=teal!4,colframe=teal!35!black,title={Before you move on}]
What two pieces is \textbf{\mr{समोर}} built from? (\emph{\textbf{\mr{सम्}}} \emph{together} and \emph{\textbf{\mr{मुख}}} \emph{face}.) Why does one word cover both \emph{in front of} and \emph{opposite}? (\textbf{A face points one way; both meanings are that}.)

Sources: \href{https://en.wiktionary.org/wiki/\%E0\%A4\%B8\%E0\%A4\%AE\%E0\%A5\%8B\%E0\%A4\%B0}{Wiktionary: \mr{समोर}}.
\end{tcolorbox}
\section[lāmb]{\mr{लांब} (lāmb) — far, and long}

\label{lesson:MR-C43-laamb}



\begin{quote}
\textbf{\mr{जवळ}} answers \emph{near}. Nothing yet answers the other half of the question.
\end{quote}

\subsection*{You'll want to know: \mr{लांब}}
\begin{quote}
\textbf{\mr{लांब}} — \emph{lāmb} — \textbf{far; long}
\end{quote}

\begin{quote}
\textbf{\mr{माझं} \mr{घर} \mr{लांब} \mr{आहे}.} — \emph{mājhaṁ ghar lāmb āhe.} — \textbf{my house is far away.}
\end{quote}

\textbf{\mr{लांब}} is \textbf{\mr{जवळ}}'s opposite and does not behave like it. \textbf{\mr{जवळ}} is a postposition and wants a noun bent in front of it; \textbf{\mr{लांब}} is an ordinary word that can simply stand in the sentence, which is why the example above needs no oblique stem at all.

The second meaning is the one to hold onto, because it explains the first. Sanskrit \textbf{\mr{लम्ब}} (\emph{lamba}) means \textquotedblleft{}hanging down, extended\textquotedblright{} — so \textbf{\mr{लांब}} is \emph{long} before it is \emph{far}, and \emph{far} is only \emph{long} applied to a distance. A Marathi speaker uses the one word for a long rope and a distant town without feeling any jump.

Pair the two and you can answer a real question:

\begin{quote}
\textbf{\mr{जवळ}, \mr{की} \mr{लांब}?} — \emph{javaḷ, kī lāmb?} — \textbf{near, or far?}
\end{quote}

\subsection*{Your turn}
Say these aloud:

\begin{itemize}
\raggedright
  \item \emph{lāmb} — far, and long
  \item \emph{mājhaṁ ghar lāmb āhe}
  \item \emph{javaḷ, kī lāmb?} — and answer it about where you live
\end{itemize}

\begin{tcolorbox}[breakable,colback=teal!4,colframe=teal!35!black,title={Before you move on}]
Which of \textbf{\mr{जवळ}} and \textbf{\mr{लांब}} needs the oblique stem? (\textbf{\mr{जवळ}} — it is a postposition; \textbf{\mr{लांब}} is not.) What does \textbf{\mr{लांब}} mean before it means \emph{far}? (\textbf{Long}.)

Sources: \href{https://en.wiktionary.org/wiki/\%E0\%A4\%B2\%E0\%A4\%BE\%E0\%A4\%82\%E0\%A4\%AC}{Wiktionary: \mr{लांब}}.
\end{tcolorbox}
\section[javaḷ · samor · lāmb]{Give somebody a direction}

\label{lesson:MR-R43-how-far}



\begin{quote}
Three words landed in this chapter and only two of them are postpositions. Name all three, then say which one is the odd one out and why.
\end{quote}

\subsection*{Your turn}
\begin{itemize}
\raggedright
  \item \emph{Say it:} \emph{ghar kuṭhe āhe?} — then answer \emph{puṇyājavaḷ}
  \item \emph{Say it:} \emph{khoolīsamor} — in front of the room
  \item \emph{Say it:} \emph{lāmb} — and notice you added no ending to anything
  \item \emph{Recall:} ask \emph{kadhī?} and \emph{kitī?} about the same house
\end{itemize}

\begin{tcolorbox}[breakable,colback=teal!4,colframe=teal!35!black,title={Before you move on}]
You can now answer \textbf{\mr{कुठे}?} five ways without naming a street. Which of this chapter's three words is not a postposition? (\textbf{\mr{लांब}}.)
\end{tcolorbox}
